% !TEX program = xelatex
%
% 필요 엔진: XeLaTeX 또는 LuaLaTeX (kotex 사용, pdfLaTeX 로는 조판되지 않음)
%   컴파일:  xelatex -interaction=nonstopmode analysis/paper-ko.tex
% 파일 인코딩: UTF-8
%
\documentclass[11pt]{article}

\usepackage{kotex}
\usepackage[a4paper,margin=28mm]{geometry}
\usepackage{booktabs}
\usepackage{array}
\usepackage{caption}
\usepackage{microtype}
\usepackage[hidelinks]{hyperref}

\captionsetup{font=small,labelfont=bf,skip=6pt}
\setlength{\parskip}{0.45em}
\setlength{\parindent}{0pt}
\renewcommand{\arraystretch}{1.15}

\newcommand{\floor}{\emph{\small 유의: 이 논문의 모든 중앙값은 \textbf{이미 성공한
사람들} 사이에서 측정한 도달 기간이다. 표본은 결과를 조건으로 추출되었고 분모가
없다. 따라서 여기의 어떤 수치도 성공 확률에 대한 진술의 근거가 될 수 없다.}\par\medskip}

\title{\bfseries 첫 성과까지 얼마나 걸리는가\\[0.35em]
\large 학업 종료부터 첫 1{,}000만 달러(불변가격) 돌파까지의 기간 추정,
그리고 그 편향의 목록}
\author{\normalsize 연구 저장소: \texttt{\textasciitilde/trajectories}}
\date{\normalsize 2026년 8월 23일}

\renewcommand{\abstractname}{요약}

\begin{document}
\maketitle

\begin{abstract}
\noindent 열거된 제3자 명단에서만 추출한 235명의 경력 시점 자료를 수집하고, 학업을
마친 시점부터 본인의 사업체가 \textbf{2026년 불변가격 기준 1{,}000만 달러}를 처음
넘어선 해까지의 기간을 측정했다. 1차 매출 근거로 이 기간을 계산할 수 있는 93명에서
중앙값은 \textbf{9.0년}, 95\% 부트스트랩 신뢰구간은 \textbf{[8.0, 10.0]}이다. 수집은
중앙값 안정성과 신뢰구간 폭에 대해 사전에 정한 규칙에 따라 종료했다.

이 추정치는 결국 돌파에 성공한 사람만을 기술한다. 수십 년을 일하고도 끝내 넘지 못한
사람들은 설계상 표본에서 빠져 있으며, 측정된 두 개의 가장 강한 편향은 모두 이 수치를
\textbf{짧아지는} 방향으로 민다. \textbf{이 숫자는 예측이 아니라 하한선이다.} 가장
쓸모 있는 결과는 통합 중앙값이 아니라 시대별 분할이다. 1995년 이전에 돌파한 경력은
중앙값 13.0년, 이후는 8.0년이며, 표본의 26\%만이 1995년 이전이다.
\end{abstract}

\floor

\section{무엇을 측정했는가}

각 인물에 대해 출처 URL과 신뢰도 등급을 갖춘 여섯 개의 시점을 기록했다: 출생,
학업 종료(\emph{마지막으로 취득을 완료한 학위}\,---\,중퇴자는 중퇴 연도가 아니라 그
이전에 완료한 학위 연도로 기록), 성과가 나온 분야에서의 첫 유급 노동, 첫 독립 사업,
\textbf{첫 성과}, 그리고 훨씬 뒤의 대규모 이정표.

주 지표는 학업 종료에서 첫 성과까지다. 만 18세 기준과 첫 창업 기준의 두 대안 지표는
결측 구조가 서로 달라 살아남는 표본이 다르므로 \ref{sec:robust}절에서 함께 보고한다.

\subsection{기준액은 물가로 조정되며, 이는 들리는 것보다 훨씬 중요하다}

성과는 \textbf{2026년 불변가격 기준} 연매출 1{,}000만 달러다. 명목 1{,}000만 달러는
시대를 가로질러 동일한 시험이 아니다. 1960년에 1{,}000만 달러를 넘으려면 실질 기준으로
오늘날의 약 11배 규모가 필요했다. 1960년의 기준선은 \textbf{919{,}417달러}다. 설계상
26\%가 1995년 이전인 표본에 명목 기준을 적용했다면, 이 표집틀이 가장 애써 포함시킨
바로 그 경력들을 부풀렸을 것이다.

이는 가상의 우려가 아니었다. 초기 예비조사에서 1960년대 기업에 명목 1{,}000만 달러를
적용한 탓에 한 행의 연도가 11년 늦게 기록된 적이 있다. 최종 자료의 모든 매출 연도는
해당 연도의 계산된 기준선과 대조했고, 각 행의 비고에는 사용한 연도의 환산 도구 출력이
그대로 붙어 있다.

\subsection{표본 선정}

이름은 오직 지정된 제3자 명단을 발표된 순서대로\,---\,순위표는 순위 순으로, 수상은
연도 순으로\,---\,열거하여 나오는 대로 취하는 방식으로만 진입했다. 국가, 시대, 성별은
기록하되 선정에는 결코 사용하지 않았다. 진입의 유일한 기준은 금액 요건이다.

\section{결과}

\begin{table}[h]
\centering
\begin{tabular}{lr}
\toprule
조사 인원 & 235 \\
포함 & 201 \\
제외 & 34 (14\%) \\
1차 매출 근거로 포함 & 121 \\
\textbf{주 지표 계산 가능} & \textbf{93} \\
\midrule
\textbf{학업 종료--첫 성과 중앙값} & \textbf{9.0년} \\
95\% 부트스트랩 신뢰구간 & [8.0, 10.0] \\
신뢰구간 반폭 & 1.00 \\
\bottomrule
\end{tabular}
\caption{주요 결과. 신뢰구간은 10{,}000회 재표집, 고정 시드의 백분위 부트스트랩이다.}
\end{table}

수집 순서에 따른 웨이브별 중앙값은
10.0, 9.0, 8.8, 9.0, 9.0, 9.8, 9.2, 9.0, 9.0이다.
추정치는 대략 40번째 인물 이후로 9년 부근에서 안정적이었으며, 9.8로 한 번 벗어난 것은
바로 다음 웨이브에서 되돌아왔다.

\subsection{시대별 분할이 더 쓸모 있는 결과다}

\begin{table}[h]
\centering
\begin{tabular}{lrrr}
\toprule
성과 시대 & $n$ & 중앙값 & 95\% 신뢰구간 \\
\midrule
1995년 이전 & 33 & 13.0년 & [9.0, 18.0] \\
1995년 이후 & 60 &  8.0년 & [6.5, 9.0] \\
\bottomrule
\end{tabular}
\caption{통합 중앙값은 상수가 아니라 혼합물이다.}
\end{table}

통합값 9.0은 야심 찬 경력이 지닌 안정적 속성이 아니다. 5년 차이가 나는 두 시대의
가중평균이며, 표본의 74\%가 1995년 이후다. 오늘의 경력에 이 숫자를 적용하려는 사람은
8.0을 읽어야 하고, 역사적으로 적용하려는 사람은 13.0을 읽어야 한다. 통합 수치만
보고하는 것은 이 자료가 허용하는 가장 오해를 부르는 행위다.

기전을 주장하지는 않는다. 더 빠른 자본 형성, 소프트웨어 유통의 경제성, 그리고 최근
비상장 매출의 더 나은 공시\,---\,이 셋 모두 같은 양상을 만들어내며, 마지막 것은 세상의
변화가 아니라 측정의 인공물이다.

\section{강건성}
\label{sec:robust}

\subsection{학업 종료 연도 결측이 적격 표본의 23\%를 제거한다}

1차 매출 근거로 성과 연도가 확정된 121명 중 \textbf{28명은 출처 있는 학업 종료
연도가 없어} 주 지표에 들어갈 수 없다. 이것이 추정치에 대한 가장 큰 단일 위협이며
구조적이다. 주 지표는 학위를 기준점으로 삼는데, 독학자와 비전통적 경로를 밟은 창업자에게
바로 그 문서화된 학위가 없기 때문이다.

제외된 행은 특정 지역에 몰려 있지 않다. 28명 중 13명이 미국 기반으로, 잔류 행의
48\%와 비슷하다. 따라서 이 결측이 국가의 대리변수로 보이지는 않는다. 그러나 학력과
관련해 무작위가 아닌 것은 분명하며, 중앙값에 미치는 방향은 무해하다기보다 미지다.

\begin{table}[h]
\centering
\begin{tabular}{lrrr}
\toprule
지표 & $n$ (121 중) & 중앙값 & 95\% 신뢰구간 \\
\midrule
학업 종료 기준 & 93 (77\%) & 9.0년 & [8.0, 10.0] \\
만 18세 기준 & 105 (87\%) & 15.0년 & [14.0, 16.0] \\
첫 창업 기준 & 118 (98\%) & 4.2년 & [4.0, 5.0] \\
\bottomrule
\end{tabular}
\caption{두 대안 지표는 표본을 더 많이 보존한다. 창업 기준 지표는 98\% 포괄률과
0.50의 반폭으로, 이 연구에서 가장 잘 측정된 양이다.}
\end{table}

포괄률을 근거로 가장 강하게 옹호할 수 있는 것은 창업 기준 지표다. \textbf{결국 성공하게
될 그 일을 시작한 뒤 불변가격 1{,}000만 달러까지의 중앙값은 4.2년}이며, 이 수치는
77\%가 아니라 적격 행의 98\%에 근거한다.

\subsection{중단 규칙은 경계선에서 발동했다}

수집은 마지막 웨이브들 이전에 고정한 규칙에 따라 종료했다. 최근 두 웨이브에 걸친 중앙값
변동 0.50년 미만, 그리고 신뢰구간 반폭 1.0년 이하다. 최종 변동은 (0.20, 0.00)이었다.
최종 반폭은 \textbf{정확히 1.00}으로, 1.0을 초과할 때만 실패하는 시험을 통과했다.

이 값이 불편할 만큼 경계에 가깝기에 부트스트랩을 300개 시드에 대해 다시 돌렸다. 규칙은
\textbf{300개 중 272개(90.7\%)}에서 통과한다. 나머지 28개는 [8.0, 11.0], 반폭 1.50을
주어 중단하지 않았을 것이다. 시드는 사전에 코드에 고정되어 있었으므로 사전 규정된 절차
아래에서 이 중단은 정당하지만, 경계선 중단이며 그렇게 보고한다.

경계선 중단에 대한 자연스러운 비판은 선택적 중단\,---\,수치가 마음에 들 때까지
수집을 계속했다는 것\,---\,이다. 이에 대해, 수집 종료 하루 전 160명, 매출 기준
$n=60$ 시점에 커밋된 중간 보고서가 이미 중앙값 \textbf{9.2년}을 보고하고 있다
(\texttt{analysis/archive/}). 매출 기준 표본이 절반 넘게 늘어나는 동안 추정치는
0.2년 움직였고, 웨이브별 이력도 대략 40번째 인물 이후 같은 안정성을 보인다. 중단
규칙은 수집이 \emph{언제} 끝날지를 정했을 뿐, 수치를 정하지 않았다. 더 엄격한 규칙을
선호하는 독자는 수집이 $n=93$에서 임의로 끝났고 신뢰구간은 [8.0, 11.0]보다 좁지 않다고
간주해야 한다.

\section{편향 목록}

수집 과정에서 24개의 편향을 목록화했고, 그중 여럿은 이미 자료를 오염시킨 뒤에야
발견되었다. 수치에 영향을 주는 것들은 다음과 같다.

\textbf{생존 편향, 그리고 이것이 다른 모든 것을 압도한다.} 표본은 돌파를 조건으로
추출되었다. 20년을 일하고 끝내 넘지 못한 사람들은 드물어서가 아니라 설계상 부재한다.
이는 이 설계 안에서 교정할 수 없으며, 어떤 확률 진술도 할 수 없는 이유다.

\textbf{문서화 편향, 짧아지는 방향.} 돌파는 매출이 일찍 공시되었을 때 연도를 특정할 수
있다. 일찍 공시한 기업은 더 빠른 성장과 더 나은 자금 조달 쪽으로 치우친다. 연도를
특정할 수 없는 돌파는 늦은 연도가 아니라 제외 처리되었고\,---\,제외 34건 중 30건이
연도 미특정이다\,---\,따라서 느리고 조용한 성공은 꼬리를 늘리는 대신 표본에서 빠진다.

\textbf{회상 편향, 연구 중간에 교정.} 초기 명단 작성 에이전트가 이름을 먼저 떠올린 뒤
그에 맞는 인용을 찾았다. 그것이 만든 모든 행에 유효한 인용이 있었지만 표본은 여전히
편향되어 있었다. 회상이 먼저였기 때문이다. 이후 웨이브는 명단을 열거하고 발표된 순서대로
읽어 내려간다.

\textbf{유명세 회피 조종, 부호만 뒤집힌 같은 오류.} 떠올릴 법한 이름을 피하라는 지시를
받은 뒤의 에이전트는 열거된 명단에서 피터 틸, 비노드 코슬라, 젠슨 황을 ``너무 뻔하다''는
이유로 배제했다. 이는 부호만 뒤집힌 동일한 편향이다. 규칙은 저명도가 어느 방향으로도
기준이 되지 않는다는 것이며, 마지막 웨이브는 필 나이트, 젠슨 황, 존 홉필드를 순위대로
취하고 건너뛰고 싶었던 유혹을 기록으로 남겼다.

\textbf{상관된 관측치.} 한 회사의 공동창업자들은 하나의 돌파를 공유하는데 부트스트랩은
모든 행을 독립 추출로 취급하므로, 공동창업자 쌍은 정보를 더하지 않은 채 신뢰구간을
좁힌다. 성과 \emph{사건}(주체와 연도) 하나당 한 행이라는 규칙을 표집틀에 추가했고, 첫
적용에서 다섯 후보를 걸러냈다. 규칙 도입 이전의 쌍 하나가 1995년 Palm Computing에
남아 있으나, 두 행 모두 대체 근거이므로 주 지표 집합 밖에 있어 보고된 신뢰구간에는
영향을 주지 않는다.

\textbf{웨이브 1은 다른 설계로 수집되었다.} 금액 단일 기준을 채택하기 전이라 구성이
의도적으로 조종되었다. 표본은 두 표집 설계의 혼합이며 그 불연속은 되돌릴 수 없다.

\section{감사 방법}

각 웨이브는 추출된 표본을 \textbf{블라인드}로 재조사하여 감사했다. 1차 조사의 답을 볼 수
없는 2차 조사를, 지침서와 물가 환산 도구만 있는 별도 디렉터리에서 수행했다. 초기 감사
하나가 공용 디렉터리에서 1차 조사의 캐시 파일을 읽고 있었던 것이 발견된 뒤 도입한
조치다.

불일치는 \emph{모순}\,---\,1년을 초과해 떨어져 있고 겹치지도 않는 두 판독\,---\,으로
채점한다. 0.10을 넘으면 웨이브는 무효가 되고 그 이름들은 재사용되지 않고 폐기된다.

연구 중 두 번의 규칙 변경이 있었고, 둘 다 시험을 느슨하게 만들었기에 여기 밝힌다. 첫째,
겹치는 구간 추정은 일치로 간주한다. 정직한 구간을 단일 연도로 접도록 강요하는 것이
멀쩡한 웨이브를 무효화하고 있었다. 둘째, 누락\,---\,한쪽은 연도를 찾고 다른 쪽은 찾지
못한 경우\,---\,을 불일치로 세지 않게 했다. 6개 플랫폼 블라인드 교차 검증에서 단일 조사의
누락률이 25\%에 근접한 것으로 측정되었으므로, 10개 감사 행에서 두세 건의 누락은 검색
철저성만으로도 발생한다. 이를 세는 것은 틀린 연도가 하나도 없는 웨이브를 무효화했다. 이제
누락은 구조 신호로 취급한다.

마지막 두 감사의 모순율은 \textbf{0.000}이었다.

\section{옹호하지 않을 것들}

\begin{itemize}
\item \textbf{연간반복매출(ARR)이 역년 매출로 취급되었을 수 있다.} 한 행이 ARR로
연도가 정해졌고, GAAP 매출을 쓴 블라인드 2차 조사는 겹치는 더 넓은 구간을 주어 감사가
이를 걸러내지 못했다. 다른 SaaS 행들도 같은 방식일 수 있으나 전수 점검은 하지 않았다.
\item \textbf{세 개 행은 연도를 결정하는 근거가 단일 출처}이며, 대안 판독이 해당 행의
비고에 명시되어 있다.
\item \textbf{일곱 개의 미해결 질문이 남아 있다.} 그중 하나는 창업 연도에 관한 것으로,
교정될 경우 어떤 구간 추정이 10년 상한을 넘겨 그 행이 제외로 바뀐다. 그 행은 주 지표
집합 밖에 있으므로 보고된 중앙값은 움직이지 않는다.
\item \textbf{초기 명단 35개 행}은 인물을 명시한 출처 URL을 요구하는 규칙 이전의 것이다.
\item \textbf{수집은 문서화된 지침을 따르는 언어 모델이 수행}했고 방법론은 사람이
결정했다. 지침서, 모든 배치 로그, 모든 감사 대조표, 모든 기각 기록이 저장소에 있다.
\end{itemize}

\section{결론}

결국 2026년 불변가격 1{,}000만 달러를 넘어선 무언가를 만들어낸 사람들 사이에서, 중앙값은
\textbf{학업을 마친 뒤 9.0년}이었다. 95\% 신뢰구간 [8.0, 10.0]. 다만 1995년 이전
돌파라면 13.0년, 이후라면 8.0년이며, \textbf{성공한 그 사업을 시작한 시점부터라면
4.2년}이다.

가장 방어 가능한 해석은 가장 좁은 해석이다. 이 자료는 도제 기간이 성과로 끝난 사람들의
도제 기간이 얼마나 길었는지를 측정한다. 몇 명이 그 기간을 견디고도 아무것도 얻지 못했는지는
말할 수 없으며, 열거 절차에서 살아남은 편향들은 살아남은 수치를 짧게 민다. 9년은 지망하는
이를 위한 기댓값이 아니라, 관측된 것에 대한 하한선으로 읽어야 한다.

\end{document}
